\subsection{Transpositions and the Alternating Group}

\begin{exercise}[3.5.1]
    In \exref{1.3.1} and \exref{1.3.2} you were asked to find the cycle decomposition of some permutations. Write each of these permutations as a product of transpositions. Determine which of these is an even permutation and which is an odd permutation.
\end{exercise}

\begin{sol}
    For \exref{1.3.1}, odd permutations are $\sigma$, $\sigma\tau$, $\tau\sigma$, and $\tau^2\sigma$, while even permutations are $\tau$ and $\sigma^2$. For \exref{1.3.2}, odd permutations are $\tau$, $\tau\sigma$, and $\tau^2\sigma$, while even permutations are $\sigma, \sigma^2$, and $\sigma\tau$.
\end{sol}

\begin{exercise}[3.5.2]
    Prove that $\sigma^2$ is an even permutation for every permutation $\sigma$.
\end{exercise}

\begin{sol}
    Let $\sigma$ be a permutation with a transposition product $\tau_1\tau_2 \cdots \tau_k$. Then $\sigma^2$ has $2k$ transpositions in its product, hence is even.
\end{sol}

\begin{exercise}[3.5.3]
    Prove that $S_n$ is generated by $\set{(i\ i+1)\mid 1 \le i \le n-1}$. [Consider conjugates, viz.\ $(2\ 3)(1\ 2)(2\ 3)^{-1}$.]
\end{exercise}

\begin{sol}
    Inducting on $n$, note that the base case is trivial. Let $T_n = \gen{(i\ i + 1) \mid 1 \leq i \leq n - 1}$. Suppose that $S_n = T_n$ for some $n \geq 1$, and consider some $\sigma \in S_{n + 1}$. Write $\sigma$ as a product of transpositions, and let $(j\ k)$ be a transposition in $\sigma$.

    There are 2 cases to consider. If $1 \leq j < k \leq n$, then $(j\ k) \in S_n = T_n$ by the inductive hypothesis. If $k = n + 1$ and $j = n$, then $(j\ k) \in T_{n + 1}$ by definition. Suppose now that $1 \leq j < k = n + 1$ where $j < n$. Note that
    \[
    	(j\ n + 1) = (n\ n + 1)(j\ n)(n\ n + 1)
    \]
    where each transposition on the right-hand side is in $T_{n + 1}$ (in particular, $(j\ n) \in S_n = T_n$ by the inductive hypothesis). It follows that $(j\ k) \in T_{n + 1}$ for all transpositions in $\sigma$, hence $\sigma \in T_{n + 1}$. By induction, the result holds for all $n \geq 1$.
\end{sol}

\begin{exercise}[3.5.4]
    Show that $S_n=\langle (1\ 2),(1\ 2\ 3 \ldots n)\rangle$ for all $n\ge2$.
\end{exercise}

\begin{sol}
    Note that $(1\ 2\ 3 \ldots n)(1\ 2)(1\ 2\ 3 \ldots n)^{-1} = (2\ 3)$. We may conjugate $(2\ 3)$ by the $n$-cycle again to get $(3\ 4)$, and so on until an arbitrary $(i\ i + 1) = (1\ 2\ 3 \ldots n)^{i - 1}(1\ 2)(1\ 2\ 3 \ldots n)^{-(i - 1)}$ is obtained for each $1 \leq i < n$. By the previous exercise, then $S_n = \gen{(1\ 2), (1\ 2\ 3 \ldots n)}$.
\end{sol}

\begin{exercise}[3.5.5]
    Show that if $p$ is prime, $S_p=\langle \sigma,\tau\rangle$ where $\sigma$ is any transposition and $\tau$ is any $p$-cycle.
\end{exercise}

\begin{sol}
    It is clear that the set of all transpositions is equal to $S_n$, since any permutation in $S_n$ can be written as a product of transpositions. Let $\sigma = (i\ j)$ and $\tau$ be a $p$-cycle. Note that there exists some power $k$ such that $\tau^k(i) = 1$. Then $\sigma' = \tau^k\sigma\tau^{-k} = (1\ \tau^k(j))$ is a transposition in $\gen{\sigma, \tau}$. We claim that every transposition of the form $(1\ k) \in \gen{\sigma, \tau}$ for $2 \leq k \leq p$, and we proceed by induction. The base case of $(1\ \tau^k(j))$ is done. Suppose $(1\ m) \in \gen{\sigma, \tau}$ for some $2 \leq m < p$. Note that there must exist some power $n$ such that $\tau^n(1) = m$. Then $\tau^n(1\ \tau^k(j))\tau^{-n} = (m\ \tau^{n + k}(j))$ lies in $\gen{\sigma, \tau}$. Then $(1\ \tau^{n + k}(j)) = (1\ m)(m\ \tau^{n + k}(j))(1\ m) \in \gen{\sigma, \tau}$. By induction, every transposition of the form $(1\ k)$ is in $\gen{\sigma, \tau}$. Then for any transposition $(a\ b)$, we have $(a\ b) = (1\ a)(1\ b)(1\ a) \in \gen{\sigma, \tau}$. It follows that all transpositions are in $\gen{\sigma, \tau}$, hence $S_p = \gen{\sigma, \tau}$.
\end{sol}

\begin{exercise}[3.5.6]
    Show that $\langle (1\ 3),(1\ 2\ 3\ 4)\rangle$ is a proper subgroup of $S_4$. What is the isomorphism type of this subgroup?
\end{exercise}

\begin{sol}
    Put $\tau = (1\ 3)$ and $\sigma = (1\ 2\ 3\ 4)$. Then $\sigma^4 = \tau^2 = 1$. Moreover, $\sigma\tau = (1\ 2\ 3\ 4)(1\ 3) = (1\ 4)(2\ 3)$, and $\tau\sigma\inv = (1\ 3)(4\ 3\ 2\ 1) = (1\ 4)(2\ 3)$, so that $\sigma\tau = \tau\sigma\inv$. This satisfies the same relations in $S_4$ as $r$ and $s$ do in $D_8$. We then define a surjective homomorphism $\phi : D_8 \to \gen{\sigma, \tau}$ by $\phi(r) = \sigma$ and $\phi(s) = \tau$. Then $\gen{\sigma, \tau}$ has at most 8 elements, showing that it is a proper subgroup of $S_4$. It is easy to see that $\phi$ maps to distinct elements in $\gen{\sigma, \tau}$, hence $\abs{\gen{\sigma, \tau}} = 8$ and $\phi$ is an isomorphism. Therefore, $\gen{\sigma, \tau} \cong D_8$.
\end{sol}

\begin{exercise}[3.5.7]
    Prove that the group of rigid motions of a tetrahedron is isomorphic to $A_4$. [Recall \exref{1.7.20}.]
\end{exercise}

\begin{sol}
    By \exref{1.7.20}, we know that the group of rigid motions is isomorphic to a subgroup of $S_4$. Moreover, \exref{1.2.9} shows that this group has order 12. We claim that the subgroup of $S_4$ that this group is isomorphic to is $A_4$. 

    The group of rigid motions of a tetrahedron consists of the following elements: the identity, 8 rotations about axes through a vertex and the center of the opposite face (by $120^\circ$ or $240^\circ$), and 3 rotations about axes through the midpoints of opposite edges (by $180^\circ$). The identity is even, the 8 rotations about vertices are even (since they correspond to $3$-cycles in $S_4$), and the 3 rotations about midpoints of edges are also even (since they correspond to products of two disjoint transpositions in $S_4$). Therefore, all 12 elements of the group of rigid motions are even permutations, so the group is isomorphic to a subgroup of $A_4$. Since both groups have order 12, it follows that the group of rigid motions of a tetrahedron is isomorphic to $A_4$.
\end{sol}

\begin{exercise}[3.5.8]
    Prove the lattice of subgroups of $A_4$ given in the text is correct. (By the preceding exercise and the comments following Lagrange's Theorem, $A_4$ has no subgroup of order $6$.)
\end{exercise}

\begin{sol}
    The previous exercise shows that $A_4$ is isomorphic to the group of rigid motions of a tetrahedron, where the text shows that this group cannot have a subgroup of order 6.

    We know that a subgroup of $A_4$ with order 4 must be isomorphic to either $Z_4$ or $V_4$. Since the symmetries of $A_4$ have either order 1 (the identity), order 2 (the $180^\circ$ rotations about midpoints of opposite edges), or order 3 (the $120^\circ$ and $240^\circ$ rotations about vertices), there are no elements of order 4 in $A_4$. Therefore, the unique subgroup of order 4 in $A_4$ must be isomorphic to $V_4$, namely $\gen{(1\ 2)(3\ 4), (1\ 3)(2\ 4)}$. Moreover, this is the only subgroup of order 4 in $A_4$ since any other subgroup of order 4 would necessarily contain different elements of order 2, which there are none of besides in this subgroup.

    The remaining potential subgroup orders of $A_4$ are 2 or 3, both of which are cyclic. Since the lattice shows 4 cyclic subgroups of order 3 that contain all 8 elements of order 3 in $A_4$, and 3 cyclic subgroups of order 2 that contain all 3 elements of order 2 in $A_4$, the lattice is correct.
\end{sol}

\begin{exercise}[3.5.9]
    Prove that the (unique) subgroup of order $4$ in $A_4$ is normal and is isomorphic to $V_4$.
\end{exercise}

\begin{sol}
    The above discussion shows that $A_4$ is isomorphic to $V_4$. Moreover, this subgroup is generated by all elements of order 2 in $A_4$. By \exref{3.1.26}, this subgroup is normal in $A_4$.
\end{sol}

\begin{exercise}[3.5.10]
    Find a composition series for $A_4$. Deduce that $A_4$ is solvable.
\end{exercise}

\begin{sol}
    Since the subgroup lattice is correct, we know that the subgroup of order 4 is maximal and normal in $A_4$. Moreover, it is Abelian. Therefore, we have the composition series
    \[
    	1 \nsub \gen{(1\ 2)(3\ 4)} \nsub \gen{(1\ 2)(3\ 4), (1\ 3)(2\ 4)} \nsub A_4
    \]
    where each composition factor is of prime order. By the characterization of solvable groups, $A_4$ is solvable.
\end{sol}

\begin{exercise}[3.5.11]
    Prove that $S_4$ has no subgroup isomorphic to $Q_8$.
\end{exercise}

\begin{sol}
    If such a subgroup $H$ of $S_4$ existed, then it would contain all elements of order 4 in $S_4$ (of which there are 6) since $Q_8$ has 6 elements of order 4. However, this implies that $(1\ 2\ 3\ 4)^2 = (1\ 3)(2\ 4) \in H$ and $(1\ 4\ 3\ 2)^2 = (1\ 4)(2\ 3) \in H$. Along with the identity, this shows that $H$ contains more than 8 elements, a contradiction. Therefore, no such subgroup exists.
\end{sol}

\begin{exercise}[3.5.12]
    Prove that $A_n$ contains a subgroup isomorphic to $S_{n-2}$ for each $n\ge3$.
\end{exercise}

\begin{sol}
    Recall that $A_n$ contains only even permutations, while $S_{n - 2}$ contain both even and odd permutations. Our homomorphism must then mapp odd permutations of $S_{n - 2}$ to even permutations of $A_n$, while ensuring that the original action of the permutation in $S_{n - 2}$ is preserved. Hence, we ``add'' on a transposition to odd permutations in $S_{n - 2}$ to make them even in $A_n$. We then define the map $\phi: S_{n - 2} \to A_n$ by
    \[
    	\phi(\sigma) = 
        \begin{cases}
            \sigma & \text{if $\sigma$ is even,}\\
            \sigma \circ (n - 1\ n) & \text{if $\sigma$ is odd.}
        \end{cases}
	\]
    Now suppose $\sigma, \tau \in S_{n - 2}$. There are 3 cases to consider:
    \begin{enumerate}
        \item If both $\sigma$ and $\tau$ are even, then $\phi(\sigma\tau) = \sigma\tau = \phi(\sigma)\phi(\tau)$.
        \item If both $\sigma$ and $\tau$ are odd, then $\phi(\sigma\tau) = \sigma \circ (n - 1\ n)\tau \circ (n - 1\ n) = \phi(\sigma)\phi(\tau)$.
        \item Without loss of generality, assume $\sigma$ is even and $\tau$ is odd. Then $\phi(\sigma\tau) = \sigma\tau \circ (n - 1\ n) = \phi(\sigma)\phi(\tau)$.
    \end{enumerate}
    In all cases, $\phi(\sigma\tau) = \phi(\sigma)\phi(\tau)$, so $\phi$ is a homomorphism. It is clear that $\phi$ is injective, since if $\phi(\sigma) = \phi(\tau)$, then either both $\sigma$ and $\tau$ are even or both are odd, so that $\sigma = \tau$. Therefore, $\phi$ is an injective homomorphism from $S_{n - 2}$ to $A_n$, showing that $A_n$ contains a subgroup isomorphic to $S_{n - 2}$.
\end{sol}

\begin{exercise}[3.5.13]
    Prove that every element of order $2$ in $A_n$ is the square of an element of order $4$ in $S_n$. [An element of order $2$ in $A_n$ is a product of $2k$ commuting transpositions.]
\end{exercise}

\begin{sol}
    By the hint, every element of order 2 in $A_n$ contains a pair of cycles of the form $(a\ b)(c\ d)$. Note that $(a\ c\ b\ d)^2 = (a\ b)(c\ d)$, and $(a\ c\ b\ d)$ is of order 4 in $S_n$. Then any element of order 2 in $A_n$ can be written as the square of an element of order 4 in $S_n$ by grouping its transpositions into pairs and applying this construction to each pair.
\end{sol}

\begin{exercise}[3.5.14]
    Prove that the subgroup of $A_4$ generated by any element of order $2$ and any element of order $3$ is all of $A_4$.
\end{exercise}

\begin{sol}
    Observing the lattice was proven correct in \exref{3.5.8}, we see that any subgroup generated by an element $\sigma$ of order 2 and an element $\tau$ of order 3. Then $\gen{\sigma}$ is a proper subgroup of $\gen{\sigma, \tau}$, but the lattice shows that $\gen{\sigma}$ is maximal in $A_4$. Therefore, $\gen{\sigma, \tau} = A_4$.
\end{sol}

\begin{exercise}[3.5.15]
    Prove that if $x$ and $y$ are distinct $3$-cycles in $S_4$ with $x \ne y^{-1}$, then the subgroup of $S_4$ generated by $x$ and $y$ is $A_4$.
\end{exercise}

\begin{sol}
    Using the same argument in the previous exercise, we note that $\gen x$ and $\gen y$ are distinct, proper subgroups of $\gen{x, y}$. Since both are maximal in $A_4$, it follows that $\gen{x, y} = A_4$.
\end{sol}

\begin{exercise}[3.5.16]
    Let $x$ and $y$ be distinct $3$-cycles in $S_5$ with $x\ne y^{-1}$.
    \begin{subexercises}
    \item Prove that if $x$ and $y$ fix a common element of $\set{1,\dots,5}$, then $\langle x,y\rangle \cong A_4$.
    \item Prove that if $x$ and $y$ do not fix a common element of $\set{1,\dots,5}$, then $\langle x,y\rangle = A_5$.
    \end{subexercises}
\end{exercise}

\begin{solenum}
    \item Consider $X = \set{1, 2, 3, 4, 5} - \set i$, where $i \in \set{1, 2, 3, 4, 5}$ is the common element being fixed by both $x$ and $y$. Then $x$ and $y$ act only on $X$. Define the map $\phi : \gen{x, y} \to S_4$ given by $\phi(\sigma) = \sigma|_X$. It is clear that $\phi$ is a homomorphism. 
    
    Now suppose $\sigma \in \ker\phi$. Then $\sigma$ fixes all elements in $X$ as well as $i$, so that $\sigma$ is the identity. Therefore, $\phi$ is injective. We now look at $\phi(\gen{x, y}) = \gen{\phi(x), \phi(y)}$. Since $\phi$ is injective, it follows that $\phi(x) \neq \phi(y)$ and $\phi(x) \neq \phi(y)\inv$ if $x \neq y$ and $x \neq y\inv$. By the previous exercise, we have distinct 3-cycles in $S_4$ such that one is not equal to the other's inverse, so that $\gen{\phi(x), \phi(y)} = A_4$. Therefore, $\gen{x, y} \cong A_4$.
    \item Put $x = (a\ b\ c)$ and $y = (a\ d\ e)$ for distinct $a, b, c, d, e \in \set{1, 2, 3, 4, 5}$. Then $xy = (a\ d\ e\ b\ c)$ is a 5-cycle, hence $\gen{x, y}$ contains a subgroup of order 5. Moreover, $yxy\inv = (b\ c\ d)$ and $xyx\inv = (b\ d\ e)$, both of which are distinct 3-cycles not equal to each other's inverses. Hence, $\gen{x, y}$ contains a subgroup isomorphic to $A_4$, which has order 12. By Langrange's Theorem then $\gen{x, y}$ must have order of at least 60, hence $\gen{x, y} = A_5$.
\end{solenum}

\begin{exercise}[3.5.17]
    If $x$ and $y$ are $3$-cycles in $S_n$, prove that $\langle x,y\rangle$ is isomorphic to $\z_3$, $A_4$, $A_5$, or $\z_3\times\z_3$.
\end{exercise}

\begin{sol}
    There are 4 cases to consider:
    \begin{enumerate}
        \item If $x$ and $y$ act on the same 3 elements, they are either the same cycle or inverses of each other. In this case, $\gen{x, y} = \gen x \cong Z_3$.
        \item If $x$ and $y$ act on 4 elements, they must fix a common element, hence $\gen{x, y} \cong A_4$ by part (a) of the previous exercise.
        \item If $x$ and $y$ act on 5 elements and do not fix a common element, then $\gen{x, y} = A_5$ by part (b) of the previous exercise.
        \item If $x$ and $y$ act on 6 elements, then they are disjoint and commute. We show that $\gen{x, y} \cong Z_3 \times Z_3$. Let $z \in Z_3$ such that $\gen z = Z_3$, and define the map $\phi : \gen{x, y} \to Z_3 \times Z_3$ by $\phi(x^a y^b) = (z^a, z^b)$ for $a, b \in \set{0, 1, 2}$. It is clear that $\phi$ is a homomorphism. Suppose $\phi(x^a y^b) = (1, 1)$. Then $z^a = 1$ and $z^b = 1$, so that $a \equiv 0 \pmod 3$ and $b \equiv 0 \pmod 3$. Therefore, $x^a y^b = 1$, showing that $\phi$ is injective. It is clear that $|\gen{x, y}| = 9 = |Z_3 \times Z_3|$, since there are 3 choices for both $a$ and $b$. Hence, $\phi$ is an isomorphism, and $\gen{x, y} \cong Z_3 \times Z_3$. \qh
    \end{enumerate}
\end{sol}